\cotefile{18}
\scanpage{18}\mspage{A - 20 -}\hfill n\textsuperscript{o} 413
% archive sheet number: 17

(resp. $U_n(\underline V)$) et ses éléments seront appelés \uline{distributions} (resp. \uline{distributions d'ordre} $\leq n$) sur $\underline V$. La restriction $\varphi_*$ à $U(\underline V)$ du transposé de $\varphi^*$ est donc un isomorphisme de modules filtrés de $U(\underline V)$ sur $S(M)$ et permet de transporter à $U(\underline V)$ la structure de coalgèbre avec forme unité de $S(M)$ : la prop. 6 montre que la structure de coalgèbre obtenue est indépendante de la carte choisie et la multiplication dans $\widehat P(\underline V)$ (qui s'identifie au dual de $U(\underline V)$) est la transposée de l'application diagonale de $U(\underline V)$. On désignera par $U_+(\underline V)$ le noyau de la forme unité : c'est aussi l'image réciproque par $\varphi_*$ du sous-module $\sum\limits_{k>\unsure{0|o}}S^k(M)$. On appellera \uline{module tangent} à $V$ et on notera $T(\underline V)$ l'intersection $U_+(\underline V)\cap U_1(\underline V)$ : c'est aussi l'image réciproque de $M=S^1(M)$ par $\varphi_*$.

Si $f\in\widehat P_+(\underline V;\underline V')$, le transposé de $f^*$ envoie $U(\underline V)$ dans $U(\underline V')$. Sa restriction à $U(\underline V)$ sera notée $f_*$. C'est une application linéaire compatible avec les filtrations, avec les structures de coalgèbres avec forme unité. La restriction de $f_*$ à $T(\underline V)$ est une application linéaire de $T(\underline V)$ dans $T(\underline V')$ appelée \uline{application linéaire tangente à} $f$ et notée $T(f)$ : si $\underline V$ et $\underline V'$ sont les noyaux formels définis par des modules $M$ et $M'$, cette définition est en accord avec celle donnée au n\textsuperscript{o} 1, à condition d'identifier $T(M)$ et $M$ (et $T(M')$ avec $M'$) ce qui est loisible.

Supposons maintenant que $\underline V'$ soit le noyau formel associé à un module libre $M$. On peut alors à $\underline u\in\widehat P_+(\underline V;M)$ associer la composée de $\underline u_*$ et de la projection canonique de $S(M)$ sur $S^1(M)=M$ : on obtient ainsi un isomorphisme de $\widehat P_+(\underline V;M)$ sur le sous-module de $\operatorname{Hom}(U(\underline V);M)$ formé des applications linéaires nulles sur $U_0(\underline V)$ (il suffit pour le voir de prendre une carte de $\underline V$ et d'utiliser la remarque 2 du n\textsuperscript{o} 4). Pour tout module séparé $N$, nous pouvons définir une application formelle

\scanpage{19}
% archive sheet number: 18
Soit $X$ un schéma algébrique lisse sur \deleted{\unsure{le}} un corps $k$ isomorphe à $\mathbb C$ \unsure{(}[de façon non précisée]. Alors
\[
 A^p(X)\xrightarrow{\unsure{P|\rho}_{\mathrm{alg}}}
 \underline{\underline H}^{2p}(X,\unsure{\Omega}_X^{\cdot})
\]
est bien défini. Le choix d'un isom $\sigma$ $k\xleftarrow{\sim}\mathbb C$ \deleted{\unsure{fixe}} permet de définir
\[
 P_\sigma^{(p)}:\quad
 A^p(X)\xrightarrow{\mathrm{can.}\,P^{\deleted{\unsure{\mathrm{an}}}}_{\mathrm{top}}\;\deleted{\unsure{P^{\mathrm{an}}}}}
 \underline{\underline H}^{2p}(X,\Omega_X^{\cdot})
\]
% slanted
\margin{$\mathrm{can}=$ De \uline{Rham} $+$ \uline{Gaga}}
qui se factorise par un sous-réseau $E_\sigma^{2i}=H^{2p}(X_\sigma,\mathbb Z)$ de $\underline{\underline H}^{2p}(X,\Omega_X^{\cdot})$. Or on a
\[
 P_\sigma^{(p)}=\left(\frac{1}{\sigma^{\deleted{\unsure{1|-1}}}(2i\pi)}\right)^p P_{\mathrm{alg}},
\]
ce qui montre que $P_\sigma^{(p)}$ dépend beaucoup de $\sigma$, \inserted{$\pi$ étant \uline{transcendant}} Dans le cas où $A^p(X)\xrightarrow{\substack{P_\sigma^{\unsure{(p)}}\\\sim}}H^{2p}(X_\sigma,\mathbb Z)$ est un isom (ex : $p=1$, $X=\mathbb P^1$) on voit de \unsure{$=$|même} que le réseau $E_\sigma^{2i}$, égal à $\frac{1}{\sigma(2i\pi)}P_{\mathrm{alg}}(A^p(X))$, dépend beaucoup de $\sigma$. \uline{NB} \unsure{Rempl.çant|Remplaçant} $\sigma$ par $\overline{\sigma}$ change $P_\sigma^{(p)}$ en $-P_\sigma^{(p)}$, donc ne change pas les $E_\sigma^i=H^{\deleted{\unsure{2i}}i}(X_\sigma,\mathbb Z)$

\scanpage{20}
% archive sheet number: 19
SYSTÈMES de RACINES \mspage{(15)}\hfill Tribu n\textsuperscript{o} 59

\uline{n\textsuperscript{o} 5} - Identifier $V'$ à son dual à la p. 195. A la p. 196 décrire plus explicitement $P(R)$ et $Q(R)$ : $Q(R)=L_{\unsure{0|o}}\cap V'$, $P(R)$ est la proj. orthogonale de $L_0$ sur $V'$, (cf. résultat ajouté au \S~6, n\textsuperscript{o} 8) ; on a $P=Q+Zv'$ où $v'$ est la projection de $(1,0,\ldots,0)$, (on a $(\ell+1)v'=(\ell,-1,\ldots,-1)$.

En bas de la p. 196, ne pas identifier tout de suite $W$ et le groupe d'autom. de $R^{\ell+1}$ d'où il peut venir, mais remonter $W$ en un groupe $G$ ; déterminer $S(R^{\ell+1})$, $S(R^{\ell+1})^G$, $S(V')$ et $S(V')^W$. Diviser par la factorielle à la l.15, p. 197.

Dans (XI) dire, tout de go, que l'autom. échange $i$ et $\ell-i+1$. Ajouter l'effet de $\operatorname{Aut}(R)$ sur $P/Q$.

\uline{n\textsuperscript{o} 6} - Prendre $\ell\geq 1$... Pas de dual à la p. 199. Dire seulement que $\unsure{\gamma|\chi}=(4\ell-2)(\ell+1)$, en référant à la formule (13)...

Dans (IX) et (X) dire que $W(R)$ se compose des applic. linéaires telles que... ; c'est seulement un produit semi-direct. Diviser par la factorielle dans (7), p. 201. Remplacer \textquotedbl{}Etant\textquotedbl{} par \textquotedbl{}Comme ils sont\textquotedbl{}, (p. 201, l.6). Ce qu'on dit dans (XI) concerne le cas $\ell\geq 2$ ; ajouter que la même conclusion est évidente pour $\ell=1$.

\uline{n\textsuperscript{o} 7} - Prendre $\ell\geq 1$. Parler de \textquotedbl{}base\textquotedbl{} d'un système de racines. Dans (V) mettre un 4 au dénominateur de la forme canonique. Répéter les deux réseaux dans (VIII). \handwritten{\unsure{--}}

\uline{n\textsuperscript{o} 8} - Prendre $\unsure{l|1}\geq 2$, et dire que c'est réductible pour 2. Dire que $R^+$ est l'ensemble des $\varepsilon_i\pm\varepsilon_j$ pour $i>j$, (en effet on a $R=R^+\cup(-R^+)$). Dans (V) on a $R=R$.

\scanpage{21}
% archive sheet number: 20
\diagram{\begin{tikzcd}[ampersand replacement=\&,column sep=large,row sep=large]
 {H^*(X,\mathbb C)} \arrow[r,"\sim"] \& {H^{\unsure{*|x}}(X,\Omega_X^\cdot)}\\
 {\underline{\underline H}^*(Y,i^!(\mathbb C))} \arrow[u,"{\unsure{?|\wr}}"'] \arrow[r,"{\substack{\mathrm{comm}\\(\text{\unsure{trivial}})\\\sim}}"] \&
 {\underline{\underline H}^*(Y,i^!(\Omega_X^\cdot))} \arrow[u]\\
 {H_{\deleted{\unsure{Y}}}^*(\deleted{\unsure{X}}Y,\mathbb C)} \arrow[u,"{\unsure{P|\rho}_{\mathrm{top}}}"] \arrow[r,"\sim","{\frac{1}{2i\pi}\,\mathrm{can}}"'] \&
 {H^{\unsure{*|\alpha}}(Y,\Omega_Y^\cdot)} \arrow[u,"{P_{\mathrm{an}}}"]
\end{tikzcd}}

qui provient de la \unsure{comm.|can.} de
\diagram{\begin{tikzcd}[ampersand replacement=\&,column sep=large,row sep=large]
 {i^!(\mathbb C)} \arrow[r,"{\substack{\mathrm{can}\\\sim}}"] \& {i^!(\Omega_X^\cdot)}\\
 {\mathbb C_Y} \arrow[u,"{P_{\mathrm{top}}}"] \arrow[r,"\sim","{\left(\frac{1}{2i\pi}\right)^p\mathrm{can}}"'] \& {\Omega_Y^\cdot} \arrow[u]
\end{tikzcd}}
lui\unsure{$=$|-même} contenu dans la formule fondamentale

\uline{N.B.} On voit que l'\uline{image directe analytique} ne \uline{transforme pas} classes entières en \uline{classes} entières : il faut un facteur $(2i\pi)^p\,\unsure{?|2}$ pour que il en soit ainsi.

\diagram{\fbox{\begin{tikzcd}[ampersand replacement=\&,column sep=small,row sep=large]
 {\substack{\text{cycles}\\\text{analytiques}\\\text{de codim }i}}\arrow[r,dash] \&
 {A^i(X)}\arrow[r,"{P_{\mathrm{top}}^{\mathbb Z}}"]\arrow[d,"{\left(\frac{1}{2i\pi}\right)^pP_{\mathrm{an}}}"] \& {H^{2i}(X,\mathbb Z)}\arrow[d,"\mathrm{can}"] \& \\
 \& {H^{2i}(X,\Omega_X^\cdot)}\arrow[r,"\sim"] \& {H^{2i}(X,\mathbb C)} \& \text{est \unsure{comm}}
\end{tikzcd}}}

\scanpage{22}\mspage{A - 18 -}\hfill n\textsuperscript{o} 413
% archive sheet number: 21

7. \uline{Noyaux formels}.

DÉFINITION 3.- \uline{Soit} $A$ \uline{une} $K$-\uline{algèbre}. \uline{On appelle carte formelle de} $A$ \uline{un couple} $\varphi=(M,\varphi^*)$ \uline{formé d'un module libre} $M$ \uline{sur} $K$ \uline{et d'un isomorphisme} $\varphi^*$ \uline{de} $\widehat P(M)$ \uline{sur l'algèbre} $A$. \uline{On dit encore que} $\varphi$ \uline{est à valeurs dans} $M$. \uline{On dit que deux cartes formelles} $\varphi_i$ $(i=1,2)$ \uline{de} $A$ \uline{à valeurs dans des modules libres} $M_i$ \uline{sont compatibles s'il existe une application formelle} $\psi_1$ \uline{de} $M_1$ \uline{dans} $M_2$ \uline{et une application formelle} $\psi_2$ \uline{de} $M_2$ \uline{dans} $M_1$ \uline{telles que} $\varphi_2^*=\psi_1^*\circ\varphi_1^*$ \uline{et} $\varphi_1^*=\psi_2^*\circ\varphi_2^*$.

Il est immédiat que la relation de compatibilité entre cartes formelles d'une même algèbre $A$ est une relation d'équivalence . Cela nous permet de poser la :

DÉFINITION 4.- \uline{On appelle noyau formel} $\underline V$ (\uline{sur l'anneau} $K$) \uline{la donnée d'une} $K$-\uline{algèbre notée} $\widehat P(\underline V)$ \uline{et d'une classe d'équivalence de cartes formelles de} $\widehat P(\underline V)$.

Naturellement la donnée d'\uline{une} carte formelle d'une algèbre $A$ définit un noyau formel.

DÉFINITION 5.- \uline{Soient} $\underline V$ \uline{et} $\underline V'$ \uline{deux noyaux formels}. \uline{Un morphisme} $f$ \uline{de} $\underline V$ \uline{dans} $\underline V'$ \uline{est la donnée d'un homomorphisme} $f^*$ \uline{de l'algèbre} $\widehat P(\underline V')$ \uline{dans l'algèbre} $\widehat P(\underline V)$ \uline{tel que pour toute carte formelle} $\varphi=(M,\varphi^*)$ \uline{de} $\underline V$ \uline{et toute carte formelle} $\psi=(N,\psi^*)$ \uline{de} $\underline V'$, \uline{l'homomorphisme} $\varphi^{*-1}\circ f^*\circ\psi^*$ \uline{de} $\widehat P(N)$ \uline{dans} $\widehat P(M)$ \uline{soit égal à l'homomorphisme} $\underline g^*$ \uline{défini par un élément} $\underline g\in\widehat P_+(M;N)$.

Il est immédiat qu'il suffit que le critère de la définition 5 soit vérifié pour un couple $(\varphi,\psi)$ de cartes formelles. Par ailleurs on vérifie aisément que les morphismes ainsi définis ne manquent d'aucune des mirifiques propriétés fonctorielles nécessaires pour avoir droit au soleil. On peut donc définir l'\uline{isomorphisme} de deux noyaux formels, ce que nous laisserons faire au prochain rédacteur, au lecteur, au Pape ou au Général.

\scanpage{23}
% archive sheet number: 22
On généralise la formule fondamentale au cas d'un \uline{morphisme} \uline{propre} \uline{quelconque}
\[
 f:X\longrightarrow Y
\]
\unsure{de} \deleted{\unsure{deux}} variétés \deleted{\unsure{analytiques|algébriques}} holomorphes, \unsure{partout|partant|purement} de codim relative $p$ :
\diagram{\fbox{\begin{tikzcd}[ampersand replacement=\&,column sep=huge,row sep=large]
 {H^{\deleted{\unsure{2i}}\unsure{i|*|2i}}(X,\mathbb C)}
 \arrow[r,"{f_*^{\mathrm{top}}}"] \arrow[d,dash,"{\wr}"'] \&
 {H^{\deleted{\unsure{2i}}\unsure{i+p|i+2p|2i-2p}}(Y,\mathbb C)} \arrow[d,dash,"{\wr}"']\\
 {H^{\deleted{\unsure{2}}i}(X,\Omega_X^\cdot)}
 \arrow[r,"{\frac{1}{(2i\pi)^p}f_*^{\mathrm{an}}}"] \&
 {H^{\deleted{\unsure{2}}i+\unsure{p|2p}}(Y,\Omega_Y^\cdot)}
\end{tikzcd}}}

D'où généralisation analogue pour des variétés \unsure{algébriques|algébriq.} sur $\mathbb C$ ....

\scanpage{24}
% archive sheet number: 23
SYSTÈMES de RACINES \mspage{(6)}\hfill Tribu n\textsuperscript{o} 59

peut pour tout entier $n\geq 1$ trouver une suite d'éléments $x_1,\ldots,x_n\in P$ telle que $y_i=x_1+\ldots+x_i\in P$, $(1\leq i\leq n)$. Mais si $n>\operatorname{Card}P$, alors il existe deux indices $i$, $j$, $(1\leq i,j\leq n,i\neq j)$, tels que $y_i=y_j$, ce qui contredit (2).

4) Pour démontrer la proposition, il suffit, en vertu du cor. 2 de la Prop. 14, de montrer qu'il existe une base $(e_k)_{1\leq k\leq\ell}$ de $V$, telle que pour la relation d'ordre lexicographique définie par cette base, tout élément de $P$ est $>0$. Procédons par récurrence sur $\ell=\dim V$. Le cas $\ell=1$ résultant de (1), on peut supposer $\ell\geq 2$. Soit $z\in V-0$ tel que $(z,x)\geq 0$ pour tout $x\in P$, qui existe d'après 3) et soient $W$ l'hyperplan orthogonal à $z$, et $V'$ le sous-espace de $W$ engendré par $R\cap W$. Alors $R\cap W$ est un système de racines dans $V'$ et $P\cap W$ est clos dans $R\cap W$. Par hypothèse de récurrence, il existe une base $(f_i)_{1\leq i\leq\ell'=\dim V'}$ de $V'$ telle que les éléments de $P\cap W$ soient $>0$ pour l'ordre lexicographique défini par cette base. Il est alors clair que la base $(e_i)_{(1\leq i\leq\ell)}$ où $e_1=z$, $e_{i+1}=f_i$, $(i=1,\ldots,e')$, et où $e_j(j>e'+1)$ est dans le supplément orthogonal de $V'$ dans $W$ a les propriétés requises.

\uline{n\textsuperscript{o} 7}, (plus grande racine, ancien n\textsuperscript{o} 6).

Le présent Congrès n'en dit rien.

\uline{n\textsuperscript{o} 8}, (pouahs et copouahs).

Le réseau engendré par $R$ se notera $Q(R)$, et ses éléments s'appelleront les pois radiciels ; idem $Q(R^*)$ ; le réseau associé à $Q(R^*)$ sera noté $P(R)$, et ses éléments seront les poix ; idem $P(R^*)$. Donner les inclusions.

\scanpage{25}
% archive sheet number: 24
$\mathbb C$, on a un objet canonique \deleted{de} $\unsure{\xi|\zeta}$ de $T_{\mathbb C}(\mathbb Z)$ [$T$ le torseur fondamental associé à $\mathbb C$).
\[
 H^i(X,\mathbb Q(0))\rightsquigarrow H^i(X_{\mathrm{top}},\mathbb Z)
\]
L'hom $H^i(X^n,\mathbb Q(0))\longrightarrow H^{i}(Y^{\unsure{m|''}},\mathbb Q(0))$ défini par \deleted{$\unsure{A^n}$} $z\in A^n(X\times Y)$ est celui défini, \inserted{par voie \unsure{transc.}, à l'aide de} \deleted{par} la classe $\unsure{P|\rho}_{\mathrm{top}}(z)\left[=\frac{1}{(2i\pi)^n}P_{\mathrm{alg}}(z)\right]$
\[
 \boxed{\xi(z)(x)=\mathrm{pr}_{2*}^{\mathrm{top}}\bigl(\mathrm{pr}_1^*(x)P_{\mathrm{top}}(z)\bigr)}
\]
Considérons l'\deleted{\unsure{isom}} can (pour $X,i$ donnés)
\diagram{\begin{tikzcd}[ampersand replacement=\&,column sep=small,row sep=large]
 {\xi_{\mathbb C}(H^{i}(X,\mathbb Q(0))=\xi(H^i(X,\mathbb Q(0))\otimes_{\mathbb Z}\mathbb C}\arrow[r,"\sim"] \&
 {\underline{\underline H}^{\deleted{\unsure{*}}i}(X,\unsure{\Omega|\underline{\Omega}}_X^*)=\xi_{\mathbb C,\infty}(H^i(X,\mathbb Q(0))}\\
 {H^{i}(X_{\unsure{\mathrm{top}}},\mathbb C)}\arrow[u,phantom,"{\shortparallel}" description]\arrow[ur,dash,"{\sim\;\deleted{\unsure{\mathrm{can}|\mathrm{car}}}}","{\substack{\mathrm{can\ par}\\\text{Hodge-Gaga}}}"'] \&
\end{tikzcd}}

Cet isom \uline{\deleted{\illegible{1}} est \deleted{\illegible{1}} fonctoriel} en $H^{\unsure{p|i}}(X,\mathbb Q(0))$.

En effet, \unsure{si} $z\in A^n(X\times Y)$, on a \deleted{Pour le rendre canonique, il faut}
\[
\begin{aligned}
 x\ \rightsquigarrow\ \mathrm{pr}_{2*}^{\mathrm{top}}\bigl(\mathrm{pr}_1^*(x)P_{\mathrm{top}}(z)\bigr)
 &=\mathrm{pr}_{2*}^{\mathrm{top}}\left(\mathrm{pr}_1^{\unsure{a|*|\mathrm{an}}}(x)\left(\frac{1}{2i\pi}\right)^n P_{\mathrm{alg}}(z)\right)\\
 &=\left(\frac{1}{2i\pi}\right)^{-n}\mathrm{pr}_{2*}^{\mathrm{alg}}\left(\mathrm{pr}_1^*(x)\left(\frac{1}{2i\pi}\right)^n P_{\mathrm{alg}}(z)\right)
\end{aligned}
\]
\diagram{\begin{tikzcd}[ampersand replacement=\&,column sep=huge,row sep=huge]
 {H^{\unsure{p|i}}(X,\mathbb C)}\arrow[r,"{\xi_{\mathbb C}(z)\deleted{\unsure{\mathrm{alg}}}}","{\substack{\text{défini par}\\ P_{\mathrm{top}}(z)\\=(\frac{1}{2i\pi})^nP_{\mathrm{alg}}(z)}}"']\arrow[d,dash,"{\wr\;\mathrm{can}}"] \&
 {H^{\unsure{p|i'}}(Y,\mathbb C)}\arrow[d,dash,"{\wr\;\mathrm{can}}"]\\
 {\underline{\underline H}^p(X,\Omega_X^\cdot)}\arrow[r,"{\substack{\xi_{\mathbb C\infty}^{\unsure{\mathrm{de}}}(z)\\\text{défini par }P_{\mathrm{alg}}(z)\\=(\frac{1}{2i\pi})^n\xi_{\mathbb C}(z)}}"' ] \&
 {H^{\unsure{p|i'}}(Y,\underline\Omega_{\unsure{x|X|Y}}^\cdot)}
\end{tikzcd}}
% slanted
\margin{Cette compatibilité peut déjà \unsure{s'énoncer} pour analytique $\unsure{\longleftrightarrow|\longrightarrow}$ \unsure{topologique} (avec \deleted{\unsure{classes|chgt}} de corr. analytiques)}
\[
 x\ \rightsquigarrow\ \mathrm{pr}_{2*}^{\mathrm{alg}}\bigl(\mathrm{pr}_1^*(x)\underbrace{P_{\mathrm{alg}}(z)}\bigr)
\]
\hfill Ouf !

\scanpage{26}
% archive sheet number: 25
SYSTÈMES de RACINES \mspage{(12)}\hfill Tribu n\textsuperscript{o} 59

\S~9. \uuline{Classification des 6tèmes de racines}.

\uline{n\textsuperscript{o} 1} - Dire qu'on est sur $R$, (l.2 du n\textsuperscript{o} 1). A la l.5 bas de la p. 177 remplacer \textquotedbl{}Il s'agit\textquotedbl{} par \textquotedbl{}Nous allons\textquotedbl{}. En bas de la p. 177, indexer par un ensemble $I$ à $\ell$ éléments, ce qui simplifie. Remonter la prop. 1 à la p. 178, (les murs étant indexés par $I$ et $I'$). Dessiner qq. chose dans le trou de la p. 178.

\unsure{-} Il est antibourbachique que les lemmes aient des corollaires : les corollaires des pp. 179 sqq. seront qualifiés de lemmes. Dans le lemme 1, considérer l'ensemble $J$ des $j$ liés à $i$. Mieux expliciter le procédé d'identification de la p. 180 : on part d'une arête $(i,j)$ d'ordre 3 ; on enlève $i$, $j$ de l'ensemble des sommets, et on y met un nouveau sommet $p$ ; les arêtes sont, d'une part les anciennes ne contenant ni $i$ ni $j$, d'autre part, si $(i,k)$ ou $(j,k)$ est une vieille arête, $(p,k)$ sera une nouvelle arête ; les ordres sont conservés ; algébriquement, ça revient à remplacer les variables $\xi_i$ et $\xi_j$ par une nouvelle variable $\xi_p$.

Enoncé du haut de la p. 181, (devenu lemme) : ou bien un seul point de ramif. et toutes arêtes d'ordre 3, ou bien chaîne et au plus une arête d'ordre $\geq 4$.

Les lemmes 3 et 4 ont l'air, soit canulés, soit d'énoncés tirés par les cheveux ; revenir à l'état 3, (n\textsuperscript{o} 383), où c'était plus clair.

La prop. 1 remonte au début du n\textsuperscript{o}.

\uline{n\textsuperscript{o} 2} - Commencer par : nous allons voir que ces trois zèbres correspondent à des groupes ; y peuvent pas provenir de systèmes de racines à cause des 5 ou des $p\geq 7$, (référence) ; par contre, nous verrons plus loin que tous les autres proviennent de systèmes de racines. A la fin du n\textsuperscript{o}, dire que les formes obtenues

\scanpage{27}
% archive sheet number: 26
Soit $K$ un corps top. isomorphe à $\mathbb C$ [isom non choisi !]

Alors on a le foncteur \inserted{mult.} $\unsure{\xi|\zeta}=\xi_K$ sur la catégorie des motifs sur $K$, à valeurs en des $\mathbb Z$-modules, et des isom
\[
\begin{aligned}
 \xi\otimes_{\mathbb Z}\mathbb Z_{\unsure{\ell|l}} &\overset{\alpha_\ell}{\simeq} \xi_\ell\\
 \xi\otimes_{\mathbb Z}K &\overset{\alpha_\infty}{\simeq} \xi_K
\end{aligned}
\]
Considérons
\[
 \xi(H^2(\mathbb P^1,\deleted{\unsure{\mathbb Q}}\mathbb Z(0)))=\xi(\mathbb Z(-1))
\]
c'est isomorphe à $\mathbb Z$, l'isom n'étant pas choisi. Le choix d'un tel isom \inserted{\unsure{i.e.|ni} d'un des gén.''} revient à celui d'une \uline{orientation} sur $\mathbb P^1_K$ \unsure{[|(}qui est de façon naturelle une variété \deleted{analytique} réelle, isomorphe à $S^2$), ou ce qui revient au \unsure{$=$|même}, au choix de l'une des deux structures \unsure{complexes|orientées} sur $\mathbb P^1_K$, provenant du \deleted{\unsure{2}} \inserted{choix d'un} isom $K\simeq\mathbb C$, i.e. le choix d'\uline{\unsure{une|un}} des racines de $-1$ dans $K$. Comme $\pi\in\mathbb R\subset\unsure{K|\mathbb R}$ est bien déterminé comme

\scanpage{28}
% archive sheet number: 27
SYSTÈMES de RACINES \mspage{(5)}\hfill Tribu n\textsuperscript{o} 59

racine simple $v_i$ envoie $v_i$ en $-v_i$ et permute les autres racines positives. Donner un lemme sur les applic. $R\longrightarrow\text{groupes}$.

Les n\textsuperscript{o} 6 et 7 sont permutés.

\uline{n\textsuperscript{o} 6} - (Ensembles clos de racines ; ancien n\textsuperscript{o} 7).

Ajouter la prop. 4 de Serre, (n\textsuperscript{o} 400, Annexe II), et sa prop. 5 ; par contre, son cor. et le lemme 2, (p. A.11), passent en exer.

Pour la prop. 15, on prend la démonstration suivante, (fournie par Borel) :

\uline{p. 110}. \uline{Prop. 15}. Démonstration :

1) Remarquons d'abord que si $x=a.y$, $(x,y\in P)$, alors $a>0$. Sinon, en effet comme $a\neq-1$ vu l'hypothèse $P\cap-P=\emptyset$ on peut, en changeant éventuellement les rôles de $x$ et $y$, supposer $a=-1/2$, d'où $x+y=y/2\in P$ et $y/2,-y/2\in P$, ce qui est absurde. Il s'ensuit, compte tenu de la Prop. 4 du N\textsuperscript{o} 3 que si $x,y\in P$ et $(x,y)<0$, alors $x+y\in P$.

2) Montrons maintenant qu'aucune somme $y_1+\ldots+y_q$, $(q\geq 1)$, d'éléments de $P$ n'est nulle. On procède par récurrence sur $q$ ; le cas $q=2$ résultant de $P\cap(-P)\unsure{.}=\emptyset$, on peut supposer $q\geq 3$. Si $y_1+y_q=0$, alors $-y_1=y_2+\ldots+y_q$ donc $(-y_1,y_2+\ldots+y_q)>0$, d'où l'existence d'un indice $j$ $(1\leq j\leq q)$ tel que $(y_1,\unsure{y_j|\dot y_j})<0$. On a alors $y_1+y_j\in P$ et :
\[
 (y_1+y_j)+\textstyle\sum_{1\leq i\leq q,i\neq j}y_i=0,
\]
contrairement à l'hypothèse de récurrence.

3) Nous voulons montrer maintenant l'existence d'un élément $z\neq 0$ de $V$ tel que $(z,x)\geq 0$ quel que soit $x\in P$. Si ce n'est pas le cas, alors, vu (1), on

\scanpage{29}
% archive sheet number: 28
élément de $K$, on voit que le choix ainsi fait permet d'écrire
\[
 \alpha_\infty(u)=\frac{1}{2\sqrt{-1}\pi}v
\]
où $v$ est le générateur canonique de $H^1(\mathbb P^1,\Omega^1)$.

Pour les choix \unsure{différents} de topologies sur $K$ faisant de $K$ un corps valué complet \unsure{arch.}, on trouve ainsi tous les éléments possibles de $H^1(\unsure{P|\mathbb P}^1,\Omega^1)$.

Si $K$ est un corps \deleted{top.} isom à $\mathbb R$ [l'isom étant d'ailleurs \uline{unique}], alors on définit encore, par choix de $\mathbb C$ comme \deleted{\unsure{corps top.|compactification}} clôture algébrique \unsure{de|à} $\mathbb R$, un foncteur $\unsure{\xi|\zeta}=\xi_K$ à valeurs en les modules sur $\mathbb Z$, \deleted{et \unsure{un|une} \unsure{des} \unsure{isom|isom.}} \inserted{mais avec op. de $\mathbb Z/2\mathbb Z=\operatorname{Gal}(\overline{\mathbb R}/\mathbb R)$} Le sous $\deleted{\unsure{\xi}}\mathbb Z$-réseau \inserted{$\xi(M)$} que nous obtenons dans $\prod\limits_{\unsure{\ell|l}} M_\ell\times M_\infty{}_{\mathbb C}$ \inserted{\unsure{relatif|valeurs} : $\mathbb R\hookrightarrow\mathbb C$} est alors stable par ce groupe $\mathbb Z/2\mathbb Z$

\scanpage{30}
% archive sheet number: 29
SYSTÈMES de RACINES \mspage{(10)}\hfill Tribu n\textsuperscript{o} 59

un quotient sur un anneau $A$ qqconque ; ce quotient est unique car $d$ n'est pas diviseur de zéro. Pour (ii), on écrit $d=\sum\limits_x c_x Qe^x$, (cf. prop. 7) ; or, dans $d$, il y a un terme en $e^{w_{\unsure{0|o}}}$, et tous les autres sont plus petits, (pour l'ordre défini par $P^+$) ; comme les $x$ sont $\geq w_0$, il y a un seul terme dans la somme, avec coeff. 1.

Le haut de la p. 162, (3 lignes), aura déjà été dit.

\S~8. \uuline{Transformations de Coq-se-taire}.

\uline{n\textsuperscript{o} 1} - Le lemme 6 du \S~6, n\textsuperscript{o} 13 remonte ici.

\uline{n\textsuperscript{o} 2} - Dans le 2) en bas de la p. 163, prendre seulement $g=s_1$. Dans le lemme 1, on notera $L$ le monoïde, $\unsure{\underline V|\underline Y}$ l'ensemble des $n!$ mots et $C(X)$ le graphe. Dans l'énoncé, remplacer 2) par : \textquotedbl{} y'a un mot $A$ et un sommet $b$ tels que $u=Ab$ et $v=bA$ \textquotedbl{} ; ça simplifie la démonstration. Mieux expliquer comment la connexion du graphe donne la conjugaison voulue.

La fin de la démonstration, (p. 165), se simplifie ainsi : il y a deux cas ; si $a$ et $b$ sont non liés, c'est facile comme dans le texte ; si $a$ et $b$ sont liés, alors $a$ n'est lié à aucun élément de $A$, et on passe de $bAa$ à $Aba$ par les étapes $bAa$, $baA$, $aAb$, $Aba$.

\uline{n\textsuperscript{o} 3} - Commencer par un laïus introduisant les $m_i$, (avec déf. en forme des exposants), et donnant les formules (1), (2), (3), (numérotées), qui sont des trivialités. Puis une proposition, réduite à (ii) et (iii), où on supposera $G$ irréductible, ((ii) est faux avec $G$ réductible). Pour la démonstration, on

\scanpage{31}
% archive sheet number: 30
Corps de base top. $\simeq\mathbb R$ ou $\mathbb C$ (foncteurs de Hodge et \deleted{\unsure{de|De}}De Rham dans ce cas)

\scanpage{32}
% archive sheet number: 31
% letterhead
\begin{tabular}{c}
INSTITUT\\
DES HAUTES ÉTUDES SCIENTIFIQUES\\
BURES-SUR-YVETTE (S.-\&-O.)\\
928-48-53
\end{tabular}

\uline{Corps de base topologique isom à $\mathbb R$ ou $\mathbb C$}

$K$ corps topologique isom à $\mathbb C$ (il y a \unsure{donc|alors} \uline{deux} isom avec $\mathbb C$, mais on ne choisit pas). $K$ contient can. $\mathbb R$.

$M$ motif sur $K$. Alors \deleted{$\mathrm{DR}(M)$ \unsure{De Rham}} pour \unsure{tt} $X$ lisse projectif sur $K$, $X(K)$ est un espace topologique cpt (variété $K$-analytique \unsure{compact|compacte}), d'où $H^{*}(X,\mathbb Q)$, et on définit ainsi \uline{foncteur de Betti} $M\unsure{\rightsquigarrow|\sim}\mathrm{Betti}(M)$ sur les motifs sur $K$. $\mathrm{Betti}(M)\otimes_{\mathbb Q}K$ \deleted{\unsure{se|$\simeq$}} correspond à $H^*(X,K)$. Par l'argument habituel, ce dernier est isomorphe can. à la coh. de DeR., d'où
\[
 \mathrm{Betti}(M)_K\simeq\mathrm{DR}(M)\qquad \otimes\text{-isom can.}
\]
Ceci met sur \inserted{(le $K$-vecteur)} $\mathrm{DR}(M)$ une structure \uline{réelle}, dont le sous-espace réel est $\mathrm{Betti}(M)_{\mathbb R}$ ; soit $\sigma$ l'involution. La filtration canonique de $\mathrm{DR}(M)$, et $\sigma$, permettent de façon unique de splitter cette filtration, de façon à obtenir sur $\mathrm{DR}(M)$ une \uline{bigraduation} compatible avec le degré total, et redonnant la filtration can. à l'aide du premier degré. Le choix d'un isom $\unsure{\pi|\tau|\sigma|u}:K\unsure{\mathord{-}|\longrightarrow}\mathbb C$ met, \unsure{par} transport de str., une

